\section{Normalization}
\label{sec:normalization}

Before the transform, the encoder maps the 8-bit samples to a signed range centred on zero
(Annex A.1 of the specification):
\begin{equation}
  I'(m,n) = \frac{I(m,n) - M}{R}.
  \label{eq:normalize}
\end{equation}
The specification leaves $M$ and $R$ to the encoder. Encoder number two uses the image mean
and the largest excursion from it, scaled to 128:
\begin{equation}
  M = \frac{1}{XY}\sum_{m,n} I(m,n), \qquad
  R = \frac{1}{128}\max\bigl(I_{\max} - M,\; M - I_{\min}\bigr),
  \label{eq:MR}
\end{equation}
so that $I'$ lies in $[-128, 128]$ whatever the contrast of the original. This makes the
subband variances, and with them the quantizer design of Section~\ref{sec:quantization},
independent of the image's brightness and contrast.

The decoder inverts the mapping and rounds to the nearest integer in the 8-bit range:
\begin{equation}
  \hat I(m,n) = \min\Bigl(255,\ \max\bigl(0,\ \floor{\hat I'(m,n)\,R + M + \tfrac12}\bigr)\Bigr).
  \label{eq:denormalize}
\end{equation}
The specification does not spell out the rounding; this is what the NIST reference decoder
does, and the compliance tolerance of one grey level (Section~\ref{sec:compliance}) absorbs
any other reasonable choice.

Two details are easy to miss. First, $M$ and $R$ travel in the frame header as decimal
\emph{scaled integers} (Section~\ref{sec:scaled}), which keep about five significant digits.
The encoder normalizes with the exact values and the decoder denormalizes with the rounded
ones; the mismatch is a relative error below $10^{-4}$ and has no visible effect. Second,
a constant image has $R = 0$, and Equation~\eqref{eq:normalize} divides by zero. The
specification does not cover this case, and an implementation has to decide what to do
(Section~\ref{sec:implementation}).
